\section{Related work}\label{sec:related}

\paragraph{Language models and stock returns.} Text has long been used to predict returns, from dictionary
measures of the tone of news \citep{tetlock2007,tetlock2008} to sentiment scores learned from returns
\citep{ke2019}. Large language models extend this line: \citet{lopezlira2026} find that ChatGPT's reading of
headlines predicts next-day returns, and \citet{chen2022} find that LLM representations of news predict the
cross-section of returns. \citet{kelly2026} use generative AI to explain realized market returns with the
reasons given in daily market summaries, revisiting \citet{roll1988}, who found that market, industry and
public news explain only a small part of individual stock returns; we carry their explanation ladder to the
firm level. Decision models are a newer interface
\citep{zhu2026}: the model returns probabilities for typed questions instead of text, so no free
text needs parsing and every output is a forecast that can be scored.

\paragraph{Look-ahead bias and memorization.} \citet{glasserman2024} test for look-ahead bias in GPT's
sentiment readings of news by anonymizing the companies; \citet{sarkar2024} analyze look-ahead bias in
pretrained language models used for prediction; \citet{lopezlira2025} show that LLMs recall economic and
financial data from their training period, so forecasting tests inside that period are not out of sample; and
\citet{he2026} train language models only on text available at each date. In machine learning, memorization
and benchmark contamination are documented and measured directly \citep{carlini2023,golchin2024,sainz2023}. Every
decision model we use was released after our sample ended, so none can be assumed ignorant of it. We therefore
mask the states, measure how often a model still identifies the company, and pre-register a forward test on
news published after the models' release.

\paragraph{News, attention and price reactions.} Prices incorporate firm news quickly
\citep{patell1984,busse2002} but not always completely. Earnings news drifts \citep{bernard1989}; price moves
that come with news continue while moves without identifiable news reverse \citep{chan2003,savor2012};
high-frequency evidence points to underreaction to news \citep{jiang2021}; and identifying the relevant news
sharpens the link between information and volatility \citep{boudoukh2019}. Stale or recombined news is
mispriced \citep{tetlock2011,fedyk2023}, and distracted or inattentive investors underreact
\citep{hirshleifer2009,dellavigna2009}. Behavioral models explain both underreaction and overreaction
\citep{barberis1998,daniel1998}, building on the overreaction evidence of \citet{debondt1985}, while
\citet{fama1998} argues that the two kinds of anomaly are about equally frequent and fragile to the choice of
method. The split of a daily return into an overnight and an intraday part \citep{lou2019} is central here: we
ask on which side of the open the news is priced.

\paragraph{Selective prediction and calibration.} A model that can abstain trades coverage for accuracy
\citep{geifman2017}, and whether its confidence is a useful guide depends on calibration, which modern neural
networks often lack \citep{pakdamannaeini2015,guo2017}. Language models are partly calibrated about what they know
\citep{kadavath2022}, and LLM forecasts can approach the accuracy of human forecasters
\citep{halawi2024,schoenegger2024}. We measure calibration with the pre-registered expected calibration error
and use selective prediction only in the exploratory section.

\paragraph{Pre-registration and replication in finance.} Many published return predictors do not survive
multiple-testing adjustments or replication \citep{harvey2016,hou2020}, although \citet{jensen2023} find that
most factors replicate once the evidence is aggregated properly. \citet{harvey2017} calls for pre-registration
and other safeguards in financial economics, and \citet{olken2015} discusses the benefits and limits of
pre-analysis plans in economics. For trading strategies, the deflated Sharpe ratio adjusts a backtest for the
number of trials behind it \citep{bailey2014}. We pre-register hypotheses, families and corrections, log every
trial, and report every pre-registered result.

\paragraph{Our earlier work.} \citet{kargarzadeh2026} ask when news is priced in, using millions of news
articles on about three thousand US stocks, each tagged by an LLM-distilled classifier with an event type and
several attributes. They find that the reaction is concentrated before and on the day of publication, that
markets underreact to numbers and overreact to stories, and that news carries width as well as direction. This
paper extends that work in three ways. It replaces the event tags with decision models that read the masked
news and forecast the move directly; it splits the reaction at the open and at fixed decision points, so that
it can ask whether a reader of the news can trade on it; and it moves from a descriptive study to a
pre-registered, blinded test with a single unblinding.
